\documentclass[10pt,a4paper]{amsart}
\usepackage[margin=19mm]{geometry}
\usepackage{amsmath,amssymb,mathtools}
\usepackage[scr=rsfs,cal=euler]{mathalfa}
\usepackage{xfrac}
\usepackage[unicode,hidelinks,bookmarks=false]{hyperref}
\newcommand{\ie}{i.e.\ }
\newcommand{\cf}{cf.\ }
\newcommand{\resp}{respectively\ }
\newcommand{\Subsection}{\subsection}
\providecommand{\nobreakdash}{\nobreak\hbox{-}}
\setlength{\emergencystretch}{2em}
\multlinegap0pt
\usepackage{bm}
\usepackage{bbm}
\usepackage{dsfont}
\usepackage{xspace}
\usepackage{cancel}
\usepackage{mathtools}
\usepackage{amsthm}
\makeatletter
\@ifundefined{theorem}{\newtheorem{theorem}{Theorem}}{}
\makeatother
\title[Quintic surfaces with 18 cusps]{Quintic surfaces with 18 cusps}
\author{Lev Borisov, Carlos Rito}
\date{Source: arXiv:2608.12305v1 (CC BY 4.0)}
\begin{document}
\maketitle

\noindent\emph{Source and licence.} Quintic surfaces with 18 cusps. Version arXiv:2608.12305v1 (CC BY 4.0), distributed under the Creative Commons Attribution 4.0 International licence, as recorded in the arXiv licence metadata for this exact version and shown on the abstract page https://arxiv.org/abs/2608.12305v1. Full rights notice: licenses/arxiv-2608.12305-CC-BY-4.0.txt.\par\smallskip

\noindent\textit{Editorial scope.} Counted unit: the Theorem whose source label is \texttt{None} in the article (source0.tex lines 868--871, source label \texttt{None}), delivered together with the article's own complete original proof of it (source0.tex lines 873--929, one \texttt{proof} environment of 57 lines). No mathematics is written, repaired or re-derived: statement and proof are the article's own text line for line. Required context retained uncounted: none. Every definition, display and label of those lines is retained unchanged. Omitted from this package and not redistributed: 1--867, 872--872, 930--955. Nothing inside a retained excerpt is omitted or altered: every retained body is delivered exactly as the article's lines are written, so no omission receipt of any kind accompanies this package. Citations: the retained excerpts carry no citation command at all, so no bibliography entry of the article is transcribed and no reference is redistributed. Reference adaptation: none was needed and none was made; every reference inside the delivered unit resolves inside it, so no unresolved reference or citation is produced. Printed slips kept literal: no printed word, symbol, hypothesis or number of the counted statement or of its proof was altered. Layout and class: the article's own class is \texttt{article} with packages including fontenc, inputenc, amsmath, amscd, amsfonts, amssymb, amsthm, xcolor, enumitem, tikz-cd, pdfpages, hyperref; the wrapper is the lane's pinned \texttt{amsart} editorial wrapper at 10pt on A4, which restates the theorem-like environments the retained text needs and adds the article's own macro definitions that the retained text needs (none), with extra wrapper packages (bm, bbm, dsfont, xspace, cancel, mathtools). The delivered numbering is the unit's own. Assets: no retained span contains a figure environment, a graphics-inclusion command, a PDF-page inclusion, a table or a TikZ picture; no image file is copied into this package, the package declares no asset dependency and no image file is redistributed with it. Licence: version arXiv:2608.12305v1 of this article is distributed under the Creative Commons Attribution 4.0 International licence, as recorded in the arXiv licence metadata for this exact version and shown on the abstract page https://arxiv.org/abs/2608.12305v1. A full rights notice is delivered at licenses/arxiv-2608.12305-CC-BY-4.0.txt. Characters: every retained excerpt is pure ASCII, so the delivered engine, XeLaTeX with its default Latin Modern text font, reports no missing character.
\begin{theorem}
The quintic surface reconstructed above has exactly $18$ ordinary cusps and no
other singularities.
\end{theorem}

\begin{proof}
Let $K$ be the number field obtained in the reconstruction and let
\[
        S=\{F=0\}\subset\mathbb P^3_K
\]
be the resulting quintic.  We first compute, over $K$, its
singular subscheme $\Sigma$.
The computation gives
\[
        \dim\Sigma=0,\qquad \deg\Sigma=36.
\]

We next reduce the exact surface modulo $p=1031$, obtaining a surface over
the degree-$22$ extension $\mathbb F_{1031^{22}}$.
Its singular subscheme is zero-dimensional.
Its primary components have degrees
\[
        8,6,6,2,2,2,2,2,2,2,2,
\]
while the corresponding reduced components have degrees
\[
        4,3,3,1,1,1,1,1,1,1,1.
\]
In particular, the reduced singular locus consists of exactly $18$ singular points.


The last eight points in the above decomposition are directly visible from
the construction and are cusps.  We now verify the twelve cusps coming from
the Barth--Rams decomposition.  Recall that, in the decomposition
\[
        ab-c^3=zF,
\]
the equations $a=0$ and $b=0$ define the two contact cubic surfaces, while
$c=0$ is the quadric.  The Barth--Rams points are the points of the complete
intersection $V(a,b,c)$ outside the residual plane $z=0$.

What remains to be checked is that this intersection is transverse at those
points, since this is precisely the condition which gives ordinary cusps in
the Barth--Rams construction.  Intersecting the reduced primary components
of the singular locus with $V(a,b,c)$ gives degrees
\[
        4,3,3,1,1,0,0,0,0,0,0.
\]
Hence the intersection consists of $12$
reduced points.  Thus the quadric $c=0$ meets the curve $a=b=0$
transversally at these twelve points, and the local Barth--Rams equation is
of type $\mathsf A_2$.  Therefore the quintic contains $18$ ordinary cusps and
no other singularities.

Finally, an $\mathsf A_2$ singularity contributes
length $2$ to the singular subscheme, so these eighteen cusps already
contribute $36$.
Since the singular subscheme over $K$ has degree $36$,
the characteristic-zero surface contains the corresponding
eighteen ordinary cusps and no other singularities.

\end{proof}

\end{document}
